% A golden rectangle and its recursive square subdivision.
% The black curve is a logarithmic golden spiral, not circular arcs.
% Geometry reference: https://doi.org/10.1155/2019/3149602
\begin{tikzpicture}[interval diagram,x=28mm,y=28mm]
  \pgfmathsetmacro{\goldenratio}{(1+sqrt(5))/2}
  \pgfmathsetmacro{\shrink}{1/\goldenratio}
  \pgfmathsetmacro{\eyex}{\goldenratio/(1+\shrink^2)}
  \pgfmathsetmacro{\eyey}{\shrink^2/(1+\shrink^2)}

  % A clockwise quarter-turn and contraction reproduce the next square.
  \foreach \step in {0,...,5} {
    \pgfmathsetmacro{\factor}{\shrink^\step}
    \pgfmathsetmacro{\turn}{-90*\step}
    \begin{scope}[
      shift={(\eyex,\eyey)},rotate=\turn,scale=\factor,
      shift={(-\eyex,-\eyey)}
    ]
      \draw[black!35] (0,0) rectangle (1,1);
    \end{scope}
  }
  \draw[black!55] (0,0) rectangle (\goldenratio,1);

  % Contract continuously while turning clockwise around the same centre.
  \draw[domain=0:630,samples=220,variable=\angle]
    plot (
      {\eyex+\shrink^(\angle/90)*
        (-\eyex*cos(\angle)-\eyey*sin(\angle))},
      {\eyey+\shrink^(\angle/90)*
        (\eyex*sin(\angle)-\eyey*cos(\angle))}
    );

  % Only label the two largest squares; leave the inner turns uncluttered.
  \node[interval label] at (0.43,0.39) {$1$};
  \node[interval label] at (1.3,0.68) {$\varphi^{-1}$};
  \node[above=6pt,interval label] at (\goldenratio/2,1) {$\varphi$};
  \node[left=6pt,interval label] at (0,0.5) {$1$};
\end{tikzpicture}
